% Root/B/D full actual-file and complete exact primary-source audits PASS.
\begin{proposition}[Proved: rank-linear total degree on a rational double point]
\label{D:RDP-full-sheaf-degree}
Let $(A,\mathfrak m,k)$ be an algebraic henselian or complete normal
surface local ring over an algebraically closed field, and suppose that
it is a rational double point. Let $\pi:Y\to\operatorname{Spec}A$ be
its minimal resolution, with exceptional curves $E_i$ and fundamental
cycle $Z_0=\sum_i z_iE_i$. For a reflexive $A$-module $M$ of rank $r$,
let
\[
 \mathcal F=\pi^*M/(\text{torsion}).
\]
Then $\mathcal F$ is a globally generated vector bundle, and
\[
 d_i:=\deg_{E_i}(\det\mathcal F)\ge0,\qquad
 \sum_i z_i d_i\le r,\qquad \sum_i d_i\le r.             \tag{D.FS.1}
\]
In particular, for any nonnegative real weights $w_i$,
\[
 \sum_i w_i d_i\le r\max_i w_i.                          \tag{D.FS.2}
\]
All characteristics are allowed, including the bad characteristics for
the individual rational double point types.
\end{proposition}
\begin{proof}
The full-sheaf assertion for algebraic henselian rational surface
singularities, in every characteristic, is
\citet[Section 2.1 and Proposition 2.7 with its full proof,
pp.~4 and 7--8 of the author version]{C_GustavsenIle2018}.
We spell out its duality argument to include the complete case.
The torsion-free sheaf $\mathcal F$ is globally generated because it is
a quotient of $\pi^*M$. Its pushforward equals $M$: both inclusions are
checked on the puncture and then extended by reflexivity of $M$.
A global generating surjection $\mathcal O_Y^a\to\mathcal F$, together
with rationality and the vanishing of $H^2$ for a proper morphism to the
affine local surface with fibres of dimension at most one, gives
$R^1\pi_*\mathcal F=0$. Duality therefore gives
\[
 R\pi_*R\mathcal Hom_Y(\mathcal F,\omega_Y)
       \simeq R\operatorname{Hom}_A(M,\omega_A).
\]
The right side has no positive cohomology since $M$ is maximal
Cohen--Macaulay. The sheaves
$\mathcal Ext_Y^q(\mathcal F,\omega_Y)$, $q>0$, have finite support:
$Y$ is regular and the torsion-free sheaf is locally free in codimension
one and on the puncture. The two-column spectral sequence for the
displayed duality consequently forces all these positive Ext sheaves
to vanish, as well as $R^1\pi_*\mathcal Hom(\mathcal F,\omega_Y)$.
Over a regular surface the resulting depth condition says that
$\mathcal F$ is locally free. This proof uses dualizing modules and
proper duality, so applies unchanged to a complete excellent local
rational double point; it does not use ordinary absolute differentials
over its coefficient field.

For clarity, the cycle facts needed below have an elementary
intersection proof. The curves $E_i$ are smooth rational $(-2)$-curves
with connected negative-definite intersection matrix. Start with one
$E_i$ and repeatedly add an $E_j$ for which the current positive cycle
$D$ satisfies $D\cdot E_j>0$. This procedure stays below any fixed
strictly positive anti-nef exceptional cycle: if the coefficient at $j$
already equalled that of such a bounding cycle, the inequality would
be impossible. It terminates at the minimal positive anti-nef cycle
$Z_0$, whose support is the whole connected exceptional set.
At every step $D^2=-2$. Indeed, if $D\cdot E_j\ge2$, then
$(D+E_j)^2\ge0$, contradicting negative definiteness; hence the positive
intersection is one and the square stays $-2$. Thus
\[
 Z_0^2=-2,\qquad K_Y\cdot Z_0=0,
 \qquad\chi(\mathcal O_{Z_0})=1.                       \tag{D.FS.3}
\]
Rationality and $H^2$-vanishing give
$H^1(\mathcal O_{Z_0})=0$, so $H^0(\mathcal O_{Z_0})=k$.
The sequence for the effective Cartier cycle $Z_0$ now gives
$H^1(\mathcal O_Y(-Z_0))=0$.

Every element of $\mathfrak m$ vanishes on $Z_0$. To check this, for a
nonzero such element write its pullback divisor as $D+H$, where $D$ is
exceptional and $H$ has no exceptional component. The cycle $D$ is
positive on each $E_i$, and
$D\cdot E_i=-H\cdot E_i\le0$. Hence $D\ge Z_0$ by the defining
minimality of $Z_0$. Consequently
$\mathfrak m\mathcal O_Y\subset\mathcal O_Y(-Z_0)$.

Twisting a global generating surjection onto $\mathcal F$ by $-Z_0$
and again using $H^2$-vanishing shows
$H^1(\mathcal F(-Z_0))=0$. Therefore
\[
 M=H^0(Y,\mathcal F)\twoheadrightarrow
 H^0(Z_0,\mathcal F|_{Z_0}),
\]
and the target is killed by $\mathfrak m$. It has dimension at most
$\mu_A(M)$, the minimum number of generators of $M$.
A rational double point has multiplicity two. Over the infinite
residue field choose a two-generated parameter reduction $Q$ of
$\mathfrak m$. The maximal Cohen--Macaulay multiplicity formula gives
\[
 \mu_A(M)\le\operatorname{length}_A(M/QM)=2r.            \tag{D.FS.4}
\]
Equivalently, a generic line through its quadratic tangent cone gives
a parameter ideal with $\operatorname{length}(A/Q)=2$; the same rank
formula then proves (D.FS.4).

The restriction $\mathcal F|_{Z_0}$ has no $H^1$, since
$H^1(Y,\mathcal F)=0$ and the next term is $H^2(\mathcal F(-Z_0))=0$.
The vector-bundle Riemann--Roch formula on the Cartier cycle $Z_0$ and
(D.FS.3) therefore yield
\[
 r+c_1(\mathcal F)\cdot Z_0
 =\chi(\mathcal F|_{Z_0})
 =h^0(\mathcal F|_{Z_0})\le2r.
\]
Each $d_i$ is nonnegative because $\mathcal F$ is globally generated,
and every $z_i\ge1$. This proves (D.FS.1), and (D.FS.2) follows.
\end{proof}

\begin{remark}[What this estimate does and does not control]
On an $A_{N-1}$ chain one can see the estimate directly, including when
the characteristic divides $N$. The character module represented by
$(u,w^a)$ in $uv=w^N$, $1\le a<N$, has invertible full pullback on the
standard crepant charts: the local generators are $u$ before the $a$-th
component and $w^a$ afterwards. Its valuation sequence is
$c_i=\min(i,a)$, with $c_0=0,c_N=a$, so its component degrees are
$2c_i-c_{i-1}-c_{i+1}$, equal to one only at $i=a$.
The proposition avoids a classification of modules and applies also
to $D$ and $E$ configurations.

For a reflexive sheaf on a partial model with rational double points,
this bounds the total degree of its full pullback on each resolved
block by its rank. It does not assert a uniform pole bound for gluing
that full pullback to independently chosen extensions on adjacent
core curves. Such compatibility, together with an effective sparse
object presentation, remains necessary for a degree-linear hereditary
representation count.
\end{remark}
